% Generated draft: requires mathematical and source review.
\documentclass[11pt]{article}
\usepackage[margin=25mm]{geometry}
\usepackage{fontspec}
\setmainfont{Latin Modern Roman}
\usepackage{amsmath,amssymb,mathtools,unicode-math}
\setmathfont{Latin Modern Math}
\usepackage{enumitem,needspace,xurl,hyperref}
\setlength{\parindent}{0pt}
\setlength{\parskip}{5pt}
\setlength{\emergencystretch}{3em}
\newcommand{\N}{\mathbb{N}}
\newcommand{\Z}{\mathbb{Z}}
\newcommand{\Q}{\mathbb{Q}}
\newcommand{\R}{\mathbb{R}}
\newcommand{\C}{\mathbb{C}}
\newcommand{\F}{\mathbb{F}}
\newcommand{\A}{\mathbb{A}}
\newcommand{\PP}{\mathbb{P}}
\newcommand{\E}{\mathbb{E}}
\newcommand{\eps}{\varepsilon}
\newcommand{\dd}{\,\mathrm{d}}
\begin{document}
\section*{UnsolvedMath ID: 20002829}
\textbf{Source problem code:} catalogue entry
\subsection{Definitions and mathematical statement}
For a finite group $G$, a supercharacter theory is a pair $(\mathcal X,\mathcal K)$ of partitions of $\operatorname{Irr}(G)$ and $G$, with $\{1_G\}\in\mathcal K$, $|\mathcal X|=|\mathcal K|$, and each canonical supercharacter $\sigma_X=\sum_{\chi\in X}\chi(1)\chi$ constant on every member of $\mathcal K$. Order theories by refinement of their superclass partitions. For $n\geq4$, put $G=S_n$ and $N=A_n$. The four theories discussed in the supplied aid have superclass partitions
\[
\mathcal K(m_n)=\{\text{individual conjugacy classes of }S_n\},\qquad
\mathcal K(C_n)=\{K:K\text{ a conjugacy class contained in }A_n\}\cup\{S_n\setminus A_n\},
\]
\[
\mathcal K(D_n)=\{\{1\},A_n\setminus\{1\},S_n\setminus A_n\},\qquad
\mathcal K(M_n)=\{\{1\},S_n\setminus\{1\}\}.
\]
The character partition of $m_n$ consists of singletons. That of $C_n$ has singleton blocks $\{\mathbf1_n\}$ and $\{\operatorname{sgn}_n\}$, and, on the remaining irreducibles, the orbits of $\chi\mapsto\chi\operatorname{sgn}_n$. Those of $D_n$ and $M_n$ are respectively $\{\{\mathbf1_n\},\{\operatorname{sgn}_n\},\operatorname{Irr}(S_n)\setminus\{\mathbf1_n,\operatorname{sgn}_n\}\}$ and $\{\{\mathbf1_n\},\operatorname{Irr}(S_n)\setminus\{\mathbf1_n\}\}$.

The source asks whether these give exactly four infinite families, with six or seven exceptional theories, and whether the families are nested. In the aid's precise stable formulation the remaining classification assertion is
\[
\forall n>7,\qquad \operatorname{Sup}(S_n)=\{m_n,C_n,D_n,M_n\}.
\]
The aid attributes this assertion to Ladisch's open Conjecture 6.2. It reports strict refinement $m_n<C_n<D_n<M_n$ for every $n\geq4$, and totals $|\operatorname{Sup}(S_4)|=5$, $|\operatorname{Sup}(S_5)|=5$, $|\operatorname{Sup}(S_6)|=8$, $|\operatorname{Sup}(S_7)|=4$, hence six additional theories at these ranks. These finite checks do not establish the stable assertion or supply a classification at all other ranks.

For the standard embedding $S_n\subset S_{n+1}$, fixing the last letter, superclass intersection compatibility means that the nonempty sets $K\cap S_n$, for $K\in\mathcal K(T_{n+1})$, form exactly $\mathcal K(T_n)$. The aid reports this for all four families. Define integral canonical restriction to mean that every $\operatorname{Res}^{S_{n+1}}_{S_n}\sigma_X$ is a nonnegative integer linear combination of the canonical supercharacters at rank $n$. The aid's separate restriction assertion is that precisely $M_n,D_n$ satisfy this property at every transition with $n\geq4$. Writing $R_n$ for the regular character, $u_n=R_n-\mathbf1_n$ and $v_n=R_n-\mathbf1_n-\operatorname{sgn}_n$, the relevant formulas are
\[
\operatorname{Res}u_{n+1}=(n+1)u_n+n\mathbf1_n,\qquad
\operatorname{Res}v_{n+1}=(n+1)v_n+n\mathbf1_n+n\operatorname{sgn}_n.
\]
For the degree-$n$ standard irreducible $\chi^{(n,1)}$ of $S_{n+1}$, branching is $\operatorname{Res}\chi^{(n,1)}=\mathbf1_n+\chi^{(n-1,1)}$. Consequently its canonical singleton supercharacter in $m_{n+1}$ has coefficient $n/(n-1)$ on the canonical standard supercharacter in $m_n$. For $C_{n+1}$ the sign-paired standard block similarly restricts with coefficient $n/(n-1)$ on the corresponding paired block in $C_n$. Both fail integrality for every $n\geq4$. The aid's canonical-normalization restriction property and superclass compatibility are distinct from a fully specified projective system: the source's projective-system meaning of ``nested'' needs explicit group maps and compatibility conditions, which the supplied record does not provide.
\subsection{Short English statement}
The source asks for a classification of symmetric-group supercharacter theories and a compatibility relation between ranks. The supplied aid identifies four standard theories, reports six extra theories in ranks four through seven, and separates superclass compatibility from integrality of restriction with the canonical character normalization. The assertion that there are only four theories at every larger rank remains reported as open; the precise projective-system interpretation of nesting is unspecified.
\subsection{Sources}
\begin{itemize}
\item Local registry: permanent catalogue entry, exact source code catalogue entry; \texttt{aim-representation-theory-notes.json}, \texttt{notes[100]}.
\item AIM workshop, \emph{Supercharacters and combinatorial Hopf algebras}, Problem 1.17, \url{https://aimath.org/WWN/supercharacters/supercharacters.pdf}. No proposer or year is supplied.
\item Original reference: A. Hendrickson, \emph{Supercharacter theories of cyclic $p$-groups}; no publication details are supplied. The source-verification field reports checking the three-page workshop PDF, confirms ``symmetric groups'', cites adjacent Problem 1.18's four $S_n$ families, and says Problem 1.2 uses ``nested'' for a projective system. The cyclic-group reference is reported as construction and lattice background.
\item Supplied research attempt 1, registry updated 14 August 2026: partial result, machine-generated and not peer reviewed. It attributes the stable classification to Ladisch's Conjecture 6.2 without a complete citation and reports lightweight character-table checks at ranks four through seven. Canonical supercharacters are described there as Diaconis--Isaacs supercharacters; no complete reference is supplied.
\end{itemize}
\end{document}
